\chapter{结论}

工业现场与分布式基础设施中的日志诊断，本质上不是一个单纯追求模型规模或分类精度的问题，而是一个同时受到响应时限、现场算力、网络状态、数据边界、云端成本和运维责任约束的系统决策问题。面对大量重复事件与少量复杂异常并存的长尾分布，固定规则难以持续覆盖未知模式，全量上云又会把网络往返、排队时延、调用费用和弱网不可用引入实时主链。CoTriage 从这一矛盾出发，将研究重点从“选择边缘还是云端”推进到“如何根据每个事件的价值和系统当下状态，动态分配边缘与云端能力”，由此形成面向连续日志流的云边协同故障感知与决策体系。

围绕这一目标，系统建立了端、边、云职责清晰而又能够闭环协作的总体架构。端侧贴近设备、操作系统和生产服务，负责多节点持续采集、来源标识、时间统一和事件化，使原本分散、异构、缺少共同坐标的原始日志转化为可关联、可追踪的统一事件。边缘侧承担现场实时主链，以规则、模板知识、MiniLM ONNX 轻量分类器和边缘小语言模型构成由快到慢、由确定到开放的分层推理路径，使高频已知事件能够尽早结束，困难事件能够逐级获得更强语义能力。云端不再接收无差别的原始日志洪流，而是围绕少量高风险、高不确定性事件，聚合时间邻近、设备关系、模板相似性和严重度等上下文，在受控 Token 预算内生成根因候选、处置建议和结构化解释，并将结果反馈到统一事件闭环。

云边协同能否成立，关键不在于同时部署两种模型，而在于是否存在一套能够决定“何时使用哪一侧能力”的机制。CoTriage 以 RAISE 将事件风险、不确定性、模板新颖度和历史分歧等认知因素，与网络健康、队列压力、调用预算和运维策略等运行因素纳入同一决策过程，在 \texttt{LOCAL}、\texttt{CLOUD} 和 \texttt{DEFER} 三种动作之间进行事件级选择。其中，\texttt{DEFER} 并非简单的失败回退，而是把“该事件值得获得云端增强”和“当前条件允许立即发送”两个问题主动解耦：能够在本地完成的事件继续形成现场判断，暂时无法发送的高价值任务进入 Outbox 持久化等待，网络恢复后再依据风险与到期状态重放。由此，系统把实时业务连续性与深度诊断完整性从一条相互阻塞的同步链路，转化为可以分别保障、最终汇合的双阶段过程。

在工程实现上，CoTriage 已将这一方法落实为可部署、可配置、可观测的完整软件链路。采集、边缘推理、轻量模型服务、动态路由、云端网关、可靠重放、控制服务和 Web 管理端以稳定事件、Redis Streams、HTTP 接口和持久化状态协作，各组件可以独立部署和扩展。系统不仅保存最终标签，还记录推理层级、模型与配置版本、置信度、路由动作、决策原因、云端用量、Outbox 状态及反馈结果，使每次判断能够回到具体事件、具体策略和具体运行条件。管理端进一步将实时事件流、异常态势、路由状态、云端诊断、可靠传输以及规则和模型配置组织到统一观察面中，使平台既能处理故障，也能解释故障如何被处理、分析能力为何被分配到某一侧，以及处理过程是否仍处于治理边界之内。

这种架构同时回应了工业设备与生产服务监测、分布式基础设施监测两类具有明显差异的应用需求。在工业现场，系统强调就近响应、断网可用、敏感数据最小外发和明确处置建议；在 Hadoop、ZooKeeper、Linux 与 Web 服务等分布式环境中，系统强调跨来源关联、传播链辨识和集群级上下文。两类场景共享统一事件、分层推理契约、RAISE、Outbox、云端网关和管理接口，场景差异则被限制在采集适配、领域规则、模板知识、模型配置和关联策略中。因此，新增设备协议、日志来源或业务领域时，不需要重新构造整条云边协同主链，体现出“场景适配局部化、核心链路稳定化”的系统设计思想。

从系统运行与验证结果看，CoTriage 已形成贯穿采集时间、边缘推理、路由动作、云端调用、用量统计、反馈写回和可靠传输状态的连续测量链，并完成真实部署组件之间的端到端联调。运行质量分析界面能够从总体负载下钻到单事件处理阶段，自动化执行器能够在相同输入下复现不同协同策略并归档逐样本记录。这说明系统已经由概念方案推进为可运行、可操作、可追踪和可重复验证的工程原型。更重要的是，平台明确区分带真值的离线评测、受控条件下的对照实验和线上运行观测，使准确率、时延、通信量、弱网保持与多节点一致性等结论能够对应到各自的数据来源和统计边界，而不依赖孤立截图或缺少分母的百分比。

CoTriage 所探索的并不仅是一种日志分类流程，而是一种面向受限环境的分布式智能组织方式：把确定性知识和快速判断尽可能下沉到现场，把稀缺而强大的云端认知能力集中到真正需要的事件，把暂时不可用的外部能力转化为可管理、可恢复的延期任务，再以统一事件和反馈机制连接感知、决策与治理。这样，边缘不再只是云端模型的前置过滤器，云端也不再是所有数据的默认终点；二者分别承担符合自身信息范围、资源条件和时间约束的任务，并通过动态路由形成随事件和环境变化而调整的协作关系。

面向未来，随着边缘模型、云端大模型和工业知识库持续演进，这套体系可以进一步吸收更丰富的时序关联、拓扑依赖和反馈学习能力，也可以与企业现有告警、工单、知识库和设备控制平台衔接，形成从异常发现、根因诊断到处置协同的更完整闭环。但无论模型和场景如何变化，CoTriage 所确立的基本原则保持不变：现场必须拥有不依赖云端的首要判断能力，高价值复杂事件应当按需获得更强认知资源，网络波动不应使重要任务静默丢失，任何自动决策都应留下可解释、可审计、可回滚的依据。
